\subsection{Channel-erasure-tolerant digital-product graph}
\label{sec:p2-erasure-tolerant-graph}
A whole-EDF cryptographic digest cannot detect
\emph{partial} digital reuse: all 50 independently
retrieved published files have distinct full-file
SHA-256 values, even where many channels are
byte-identical. We therefore analyse the
preprocessed EDF files as channel-by-native-record
digital product arrays. For released ID $i$,
named EEG channel $c$ and aligned native EDF
storage record $r$, let $H_{icr}$ be SHA-256 of
the exact stored digital sample bytes and
$V_{icr}$ the predicate that samples are
\emph{nonconstant within that native record}.
Define
\[
M_{ij}(r)=
\sum_{c\in\mathcal C}
\mathbf{1}\{H_{icr}=H_{jcr},
V_{icr}=V_{jcr}=1\},\qquad |\mathcal C|=29.
\]
Same-position comparisons require named channel,
sample count and record geometry agreement.
The longest compatible contiguous span is
\[
R_{ij}(k)=\max\{\ell:
\exists r,\ \forall t=r,\ldots,r+\ell-1,\
M_{ij}(t)\ge k\}.
\]
This is an exact-digital-channel erasure-tolerant
content fingerprint, \emph{not} a declaration of
shared original acquisition or person.
Cryptographic digests provide practical,
not mathematically collision-free, equality evidence.

The full independent real-byte rehash and
1,225-pair census and a separately implemented
six-file, three-pair temporal topology audit
used a common version-pinned NEMAR EDF release
with no downloaded patient labels or redistributed
waveforms. For ID 4/5, $M_{45}(r)=29$ for all
320 native records. For 47/49,
$M_{47,49}(r)=28$ on 245 records, including
a contiguous 179-record run (zero-based
indices 128--306); the P4 channel has
no exact matching native time-varying block,
which demonstrates why whole-file equality
does not detect this case. For 11/15, only
the P3 channel has 171 identical nonconstant
native records; no individual record has
two or more equal EEG channels. The remaining
1,222 pairs have no exact same-channel,
same-record EEG block match in this release.

As a \emph{retrospective}, explicitly
post-outcome example, consider the rule
$R_{ij}(28)\ge 100$: at least 28 of 29
named EEG channels equal in 100 consecutive
native storage records. This gives exactly
two release-content edges, 4/5 and 47/49,
versus zero whole-file SHA-256 identity
edges and one whole-29-channel equality
edge (4/5). The other 1,223 pairs
remain unlinked under the inspected rule.
This is \emph{not} a calibrated prospective
false-positive rate. The 28- and 100-unit
thresholds were set after examining the
named released products.

\begin{table}[ht]\centering\small
\caption{Digital-product content graph
connections on all 1,225 different released
ID pairs. EDF blocks are not independent
clinical trials; graph groups are not
clinical participant identities.}
\label{tab:p2-erasure-content-graph}
\begin{tabular}{p{.57\linewidth}rl}
\toprule
Same-release product-link criterion & Edges & Public ID pairs\\
\midrule
Whole EDF SHA-256 identity & 0 & None\\
All 29 entire time-varying EEG channels
identical & 1 & 4/5\\
At least 28 nonconstant EEG channels matched
for 100 consecutive native EDF records
& 2 & 4/5, 47/49\\
\bottomrule
\end{tabular}
\end{table}

The resulting 48 connected components
are \emph{digital-product fingerprint groups};
they must \emph{not} be reported as
48 distinct original patients.
Neither raw MAT-to-EDF export lineage
nor underlying patient identity follows
from these edges. An externally validated
decision rule and measurement of how the
actual train/test split conclusions change
are still required before claiming a
new generalisable provenance algorithm.
The real-byte evidence is in
\href{https://github.com/lindicaphxag-tech/kaggle/actions/runs/38388066957}
{the full 50-file independent checksum census},
\href{https://github.com/lindicaphxag-tech/kaggle/actions/runs/38543795987}
{the nonconstant digital block topology audit}
and their machine-readable cross-artifact
comparison.